\subsection{Composition, Restriction, and Inverse Functions}

\begin{theorem} \label{thm:bijective-iff-invertible-function} \hlt{Bijective Means Invertible}\\
A function $f:X\to Y$ is bijective iff there exists a function $g:Y\to X$ s.t. $g\circ f=\operatorname{id}_X$, $f\circ g=\operatorname{id}_Y$.\\
When such a function exists, it is unique and is denoted $f^{-1}:Y\to X$.
\begin{proof}
\begin{enumerate}[label=\roman*.]
\setlength{\itemsep}{0pt}
\item[$(\Rightarrow)$] If $f$ is bijective, define $g(y)$ as the unique $x$ with $f(x)=y$. Then $g\circ f=\operatorname{id}_X$ and $f\circ g=\operatorname{id}_Y$.
\item[$(\Leftarrow)$] If such a $g$ exists, applying $g$ to $f(x_1)=f(x_2)$ gives $x_1=x_2$, and $y=f(g(y))$ gives surjectivity.
\end{enumerate}
If $g$ and $h$ are two-sided inverses, then $g(y)=g(f(h(y)))=h(y)$, so the inverse is unique.\\
\hlt{Structural Commentary}: A true inverse undoes $f$ on both sides.\\
\hlt{Key Insight}: Bijective means exactly one preimage per output.
\end{proof}
\end{theorem}

\hlt{Placeholder}: Add algebra of composition, restrictions, extensions, identities, and image/preimage laws under composition.
